%! Systems of Differential Equations — Unit 6, Lesson 6.4 — Warm-Up (Answer Key)
\documentclass[10pt]{article}
\usepackage{linalg-article}
\usepackage{linalg-key}   % pulls in -boxes; do NOT also load -boxes
\newcommand{\TallMath}[1]{$\displaystyle #1\rule[-1.4em]{0pt}{3.2em}$}
\newcommand{\uu}{\mathbf{u}}
\newcommand{\xx}{\mathbf{x}}

\begin{document}

\pageheader{Unit 6, Lesson 6.4}{Warm-Up --- Answer Key}
\namedateperiod

\vspace{0.12in}

\begin{notesbox}{Getting Ready: Three Moves We'll Use Today}
\small Today we run a system \emph{forward in time}. These three quick moves are the whole method. Answer quickly.

\vspace{4pt}
\textbf{1. Grow or decay?} A quantity with $\frac{du}{dt}=\lambda u$ is $u(t)=u(0)e^{\lambda t}$: it \emph{grows}
if $\lambda>0$ and \emph{decays} if $\lambda<0$. For each, circle grow or decay.
\[
  u(t)=4e^{3t}:\ \ \ans{grow}\ (\lambda=3>0) \qquad\qquad
  w(t)=7e^{-2t}:\ \ \ans{decay}\ (\lambda=-2<0)
\]

\vspace{4pt}
\textbf{2. Split a vector into a combination.} Find $c_1,c_2$ so that
$c_1\begin{bsmallmatrix}1\\1\end{bsmallmatrix}+c_2\begin{bsmallmatrix}1\\-1\end{bsmallmatrix}=\begin{bsmallmatrix}5\\3\end{bsmallmatrix}$.
\[
  c_1+c_2=\ans{$5$},\quad c_1-c_2=\ans{$3$}\ \Rightarrow\ c_1=\ans{$4$},\ c_2=\ans{$1$}.
\]

\vspace{4pt}
\textbf{3. Eigenvector recall.} For $A=\begin{bsmallmatrix}1&2\\2&1\end{bsmallmatrix}$, multiply
$A\begin{bsmallmatrix}1\\1\end{bsmallmatrix}$ and read off the eigenvalue.
\[
  A\begin{bmatrix}1\\1\end{bmatrix}=\ans{$\begin{bsmallmatrix}3\\3\end{bsmallmatrix}$}=\lambda\begin{bmatrix}1\\1\end{bmatrix},
  \qquad \lambda=\ans{$3$}.
\]
\end{notesbox}

\begin{teachernote}
These three moves \emph{are} today's method, in order: \textbf{(1)} the sign of $\lambda$ reading grow/decay
off $e^{\lambda t}$ --- the engine of the whole lesson; \textbf{(2)} splitting a start vector into a
combination of $\begin{bsmallmatrix}1\\1\end{bsmallmatrix},\begin{bsmallmatrix}1\\-1\end{bsmallmatrix}$ (a quick
$2\times2$ solve), the move that finds $c_1,c_2$ in $\uu(t)=c_1e^{\lambda_1 t}\xx_1+c_2e^{\lambda_2 t}\xx_2$;
and \textbf{(3)} confirming an eigenvector/eigenvalue pair, the data the method runs on. Debrief quickly, then
name the plan: ``split the start into eigenvectors, let each run at rate $e^{\lambda t}$, and add.''
\end{teachernote}

\end{document}
